\section{Implementation Challenges}


Several technical challenges arise in implementing fraud proofs for complex operations.

\subsection{Complex Opcode Verification}

Some opcodes require substantial on-chain computation:

\begin{itemize}
    \item \textbf{CALL}: Requires full contract execution simulation
    \item \textbf{CREATE}: Needs contract deployment verification
    \item \textbf{SELFDESTRUCT}: Must verify state deletion
    \item \textbf{DELEGATECALL}: Complex context switching
\end{itemize}

Solutions include:
\begin{enumerate}
    \item Pre-compiled contracts for common operations
    \item Optimistic execution with separate fraud proofs
    \item Limiting supported opcodes in rollups
    \item Hybrid on-chain/off-chain verification
\end{enumerate}

\subsection{Memory Access Patterns}

Efficient memory verification requires careful design:

\begin{equation}
\text{MemoryProof}_{\text{size}} = O(k \times \log(n))
\end{equation}

where $k$ is the number of memory accesses and $n$ is memory size.

\subsection{Gas Limit Constraints}

L1 gas limits constrain verification complexity:

\begin{table}[h]
\centering
\caption{Gas Costs for Verification Operations}
\begin{tabular}{lr}
\toprule
\textbf{Operation} & \textbf{Gas Cost} \\
\midrule
Merkle proof verification (per level) & 1,000 \\
State root computation & 5,000 \\
Simple opcode execution & 10,000 \\
Complex opcode execution & 100,000 \\
Full one-step proof & 2,000,000 \\
\bottomrule
\end{tabular}
\end{table}

\subsection{Prover-Verifier Asymmetry}

The asymmetry between proving and verification creates opportunities and challenges:

\begin{equation}
\frac{\text{Cost}_{\text{prove}}}{\text{Cost}_{\text{verify}}} = O(n)
\end{equation}

This favorable ratio enables efficient fraud proofs but requires careful protocol design to prevent exploitation.

